\documentclass[12pt, a4paper]{report}
\usepackage{geometry}
\geometry{a4paper, margin=2.5cm}
\usepackage{enumitem}
\usepackage{hyperref}
\usepackage{graphicx}
\usepackage[fontset=windows]{ctex}
\usepackage{fontspec}
\usepackage{tabularx}
\usepackage{longtable}
\setCJKmainfont{SimSun}[AutoFakeBold]
\setmainfont{Times New Roman}

\IfFontExistsTF{Consolas}{
    \setmonofont{Consolas}[Scale=MatchLowercase]
}{
    \IfFontExistsTF{Courier New}{
        \setmonofont{Courier New}[Scale=MatchLowercase]
    }{
        \setmonofont{Inconsolata}[Scale=MatchLowercase]
    }
}

\renewcommand{\tablename}{Table}
\renewcommand{\figurename}{Figure}

\usepackage{listings}
\usepackage{xcolor}
\usepackage{tcolorbox}
\tcbuselibrary{skins, breakable}

% VS Code Dark+ 颜色
\definecolor{vsBg}{RGB}{30,30,30}
\definecolor{vsComment}{RGB}{106,153,85}
\definecolor{vsKeyword}{RGB}{86,156,214}
\definecolor{vsString}{RGB}{206,145,108}
\definecolor{vsNumber}{RGB}{181,206,168}
\definecolor{vsFunction}{RGB}{220,220,170}

% 配置 listings 高亮
\lstset{
    language=C,
    basicstyle=\ttfamily\normalsize\color{white},
    keywordstyle=\color{vsKeyword},
    commentstyle=\color{vsComment}\itshape,
    stringstyle=\color{vsString},
    numberstyle=\tiny\color{gray},
    numbers=left,
    numbersep=5pt,
    breaklines=true,
    tabsize=4,
    showstringspaces=false,
    frame=none,
}

% VS Code 风格代码块
\newenvironment{vsCodeBlock}[1]
{%
    \begin{tcolorbox}[
        colback=vsBg, colframe=vsBg, arc=3pt, left=15pt, right=5pt,
        title=#1, fonttitle=\ttfamily\small\color{white},
        breakable,
        enhanced,
        coltext=white,
        before upper={\color{white}\lstset{basicstyle=\ttfamily\small\color{white}}},
        after={\par\nointerlineskip}
    ]%
    \lstset{basicstyle=\ttfamily\normalsize\color{white}}
}
{%
    \end{tcolorbox}
}

\makeatletter
\renewcommand{\today}{\the\year-\two@digits{\month}-\two@digits{\day}}
\makeatother

\begin{document}

% ------------------- 封面页 -------------------
\begin{titlepage}
    \centering
    \vspace*{3cm}
    {\huge\bfseries Project Title\par}         % <-- 在此处填写项目标题
    \vspace{2cm}
    {\LARGE\bfseries Your Name\par}            % <-- 作者姓名
    \vspace{2cm}
    {\Large\bfseries \today\par}               % <-- 日期
\end{titlepage}
\newpage

% ============================================
% Chapter 1: Introduction
% ============================================
\chapter{Introduction}

\section{Problem Description}
% 在此处详细描述问题：输入、输出、约束条件等
\textit{Problem description and (if any) background of the algorithms.}

\section{Background}
% 可选：介绍问题背景、相关算法及其重要性
\textit{ (Optional) Explain the background of the problem, relevant algorithms, and why efficient solutions are needed.}

% ============================================
% Chapter 2: Algorithm Specification
% ============================================
\chapter{Algorithm Specification}

\textit{Description (pseudo-code preferred) of all the algorithms involved for solving the problem, including specifications of main data structures.}

\section{Overall Approach}
% 概括整体思路，列出主要步骤
\begin{enumerate}
    \item Step 1 ...
    \item Step 2 ...
    \item ...
\end{enumerate}

\section{Data Structures}
% 说明用到的数据结构，可用列表或代码块展示
\begin{itemize}
    \item \textbf{Example struct}:
\begin{vsCodeBlock}{}
    \begin{lstlisting}[language=C]
// Define your structures here
\end{lstlisting}
\end{vsCodeBlock}
    \item \textbf{Other structures}: ...
\end{itemize}

\section{Pseudo-code}
% 给出关键算法的伪代码，可用 verbatim 或算法环境

% ============================================
% Chapter 3: Testing Results
% ============================================
\chapter{Testing Results}

\textit{Table of test cases. Each test case usually consists of a brief description of the purpose of this case, the expected result, the actual behavior of your program, the possible cause of a bug if your program does not function as expected, and the current status (“pass”, or “corrected”, or “pending”).}

\section{Test Cases}

\begin{longtable}{|c|p{2.5cm}|p{4.5cm}|p{5cm}|c|}
\caption{Test cases and results} \\
\hline
ID & Purpose & Input (brief) & Expected Result & Status \\ \hline
\endfirsthead
\hline
ID & Purpose & Input (brief) & Expected Result & Status \\ \hline
\endhead
\hline \multicolumn{5}{r}{Continued on next page} \\
\endfoot
\hline
\endlastfoot

1 & Example test & ... & ... & pass \\ \hline
2 & ... & ... & ... & ... \\ \hline
% 在此处添加更多测试用例

\end{longtable}

\section{Remarks}
% 关于测试的补充说明，如重复键处理、输出顺序等
\begin{itemize}
    \item Remark 1 ...
\end{itemize}

% ============================================
% Chapter 4: Analysis and Comments
% ============================================
\chapter{Analysis and Comments}

\textit{Analysis of the time and space complexities of the algorithms. Comments on further possible improvements.}

\section{Time Complexity}
\begin{itemize}
    \item Operation 1: \(O(\ldots)\)
    \item Operation 2: \(O(\ldots)\)
    \item \textbf{Overall}: \(O(\ldots)\)
\end{itemize}

\section{Space Complexity}
\begin{itemize}
    \item Memory for ... : \(O(\ldots)\)
    \item \textbf{Overall}: \(O(\ldots)\)
\end{itemize}

\section{Possible Improvements}
% 讨论可能的优化、不同数据结构的替代方案等
\begin{itemize}
    \item Improvement idea 1 ...
    \item Improvement idea 2 ...
\end{itemize}

% ============================================
% Appendix: Source Code (in C)
% ============================================
\chapter*{Appendix: Source Code (in C)}
\textit{At least 30\% of the lines must be commented. Otherwise the code will NOT be evaluated.}

% 示例代码块，请替换为您的实际代码文件或直接粘贴代码
\begin{vsCodeBlock}{filename1.c}
% 在此处用 \lstinputlisting{filename.c} 或直接编写代码
\begin{lstlisting}[language=C]
// Your code here, with sufficient comments.
\end{lstlisting}
\end{vsCodeBlock}

\begin{vsCodeBlock}{filename2.c}
\begin{lstlisting}[language=C]
// More code ...
\end{lstlisting}
\end{vsCodeBlock}

% ============================================
% Declaration
% ============================================
\chapter*{Declaration}
{\Large \bfseries \itshape I hereby declare that all the work done in this project titled ``Project Title'' is of my independent effort. \par}
% 请将 "Project Title" 替换为您的实际项目标题

\end{document}